% TikZ 图：力电综合示意
% Day 30：综合训练
\begin{tikzpicture}[scale=0.88,>=Latex]
  % 轨道在场区边界处结束；进入场区后粒子不再受轨道约束。
  \fill[blue!8] (5.2,0.3) rectangle (8.2,2.9);
  \draw[thick] (0,0) -- (3.2,0)
    .. controls (4.2,0) and (4.2,1.6) .. (5.2,1.6);
  \foreach \x in {5.6,6.2,6.8,7.4}{
    \draw[->,blue,thick] (\x,2.65) -- (\x,1.95);
  }
  \fill[red] (0.7,0.15) circle (2.5pt) node[above] {$+q$};
  \draw[->,red,thick] (0.7,0.15) -- (1.55,0.15) node[midway,above] {$v_0$};
  \fill[red] (5.2,1.6) circle (2.2pt);
  \draw[->,red,thick] (5.2,1.6) -- (6.05,1.6)
    node[midway,above] {$v$};
  \draw[red,thick,dashed,domain=5.2:8.0,samples=60]
    plot (\x,{1.6-0.10*(\x-5.2)*(\x-5.2)});
  \node at (6.7,3.2) {匀强电场区};
  \node[font=\small,align=center] at (4.1,-0.65)
    {轨道段先用能量或牛顿定律；离轨进入电场后分解运动。};
\end{tikzpicture}
